\chapter{Keeping the economic question intact across methods}
\label{in:semantics}

An analysis remains the same analysis only if its population, intervention, information and uncertainty survive each computational step. This chapter first develops that requirement for coordinates, graph classes and sets. It then examines the evidence from local custody, recovery and interface trials, before comparing restricted Python identification, Java discovery and R/native forest routes. Each route has a different scientific domain. Their common representation makes those domains visible; it does not make their guarantees interchangeable.

\section{The unit of exchange is an analysis with meaning}

A funding service, a servicing notebook and an analyst's agent may all request
an ``effect'', while referring to different populations, actions, horizons and
units. Exchanging a floating-point number is therefore insufficient. The
proposed common object contains a scientific context, an operation, input
identities, an execution disposition and a qualified result. The local
reference implementation enforces selected parts of this contract; it does not
validate the truth of every declaration. \cite{v12_in_contract,v12_in_runtime}

For a scalar intervention query, write the context schematically as
\[
 Q=(P,\mu,a,a_0,Y,H,u,s,\mathcal T,\mathcal A,G),
\]
where $P$ is the population, $\mu$ its weighting measure, $a$ and $a_0$ the
intervention and comparator, $Y$ the outcome, $H$ the horizon, $u$ the units,
$s$ the sign convention, $\mathcal T$ the information clocks, $\mathcal A$ the
assumptions and sampling declarations, and $G$ the graph or graph class.
This mathematical notation expresses the required meaning; particular schema
fields may encode only a supported subset. The original declaration must be
retained when a narrower operation representation is used.

Changing $a-a_0$ from a ten-basis-point offer to a one-percentage-point offer
changes the question, even if the same coefficient field can store both.
Replacing customer weights by balance weights changes $\mu$. Treating a
post-decision observation as available history changes $\mathcal T$. A digest
can detect that the declared context changed. It cannot establish that the
context describes the actual banking process.

The implementation binds supplied moments and sensitivity artifacts to their
scientific context. Reusing their bytes with a changed sign, clock, unit or
target is rejected unless a new, appropriate artifact is supplied. This
prevents a consumer from reinterpreting a cached numerical object through an
unannounced schema change. It does not replace a scientific argument for the
moments or covariance. \cite{v12_in_runtime}

\section{Execution, science and qualification are different states}

A successful invocation can return unsupported identification, an empty
sensitivity set, or a whole-line confidence set. These are not necessarily
transport failures. Conversely, a timeout supplies no scientific answer even
if the requested mathematical object exists. The common response separates
invocation state, scientific disposition and qualification evidence.

This separation gives a batch operator a defensible continuation rule. A
completed negative answer can be archived and used in a report. A malformed
request needs correction. An unsupported operation needs another qualified
backend or a narrower, explicitly authorized question. A contradiction in the
declared premises needs scientific review. Automatically retrying all four
events until a finite number appears would select answers according to their
convenience.

The library's supplied-DAG back-door operation illustrates the distinction.
It can attempt an adjustment expression in its bounded domain. Failure of
that attempt does not prove general nonidentification; a different graphical
functional or additional information could be relevant. The response must
describe the actual method and its limitation. Likewise, graph syntax review
checks declared structures but does not discover their causal truth.
\cite{v12_in_runtime}

\section{Sets must retain their topology}

The local representation distinguishes unknown from empty, and permits a union
of intervals with infinite endpoints. Unknown has no asserted interval list;
an empty set has an empty list. A whole-line set has one interval unbounded in
both directions. A disconnected set keeps its components. \cite{v12_in_contract}
None should be represented by an arbitrary plotting range.

Consider the stipulated effect set
\[
 \Theta=(-\infty,-2]\ \cup\ [3,\infty)
\]
and a known monetary map $V=100\theta-300$. The image is
\[
 V(\Theta)=(-\infty,-500]\ \cup\ [0,\infty).
\]
Taking the convex hull would return the entire real line and discard a real
exclusion. Clipping the tails would falsely assert boundedness. Changing to a
loss convention $L=-V$ gives $(-\infty,0]\cup[500,\infty)$; endpoint order must
reverse when multiplying by a negative number.

If $\Theta$ has a stated coverage guarantee for $\theta$, applying a known
deterministic map preserves containment on the same event. It does not create
a guarantee for an estimated conversion factor or a randomly chosen target.
When margins, discount rates or exposure weights are estimated, joint
uncertainty or an explicitly conditional analysis is needed. Serialization
cannot make uncertain financial inputs fixed.

The same care applies to graph sets. Member identities, endpoint semantics and
evidence must travel with the alternatives. A CPDAG, PAG and selected DAG are
not interchangeable edge lists. The present schema retains declared classes
and hashed membership; its validator is not a complete recognizer or confidence
construction for arbitrary temporal equivalence classes.

\section{A contract for future numerical adapters}

Each adapter should state its accepted targets, graph classes, sampling laws,
missingness rules, numerical inputs and output uncertainty. A source package
may use an ATT where the caller requested an ATE, return pointwise intervals
where simultaneous inference is required, or approximate a set with a hull.
Those are semantic mismatches even when both sides exchange valid JSON.

The adapter must either preserve the requested meaning, return a disclosed
qualified approximation, or refuse the mapping. A proposed R, Python, Java or
C++ bridge therefore needs matched fixtures at ordinary and adverse boundaries,
including disconnected sets, weak identification and unsupported models.
Cross-language execution, licenses, data copying and complete workloads remain
qualification obligations. The existing local dispatcher is a common boundary
for that work, not evidence that those bridges already exist.


\hypertarget{when-a-conversion-preserves-the-question}{%
\section{When a conversion preserves the
question}\label{when-a-conversion-preserves-the-question}}

The interchange obligation can be stated as a commutation requirement.
Let \(F_Q(D)\) be a scientific calculation on observations \(D\) for
context \(Q\). Let \(E\) encode the observations and question for an
engine, \(M\) denote that engine's computation, and \(R\) decode its
result. If \(T\) is a declared change of coordinates, an exact adapter
should satisfy

\[
 R\{M(E(D,Q))\}=T\{F_Q(D)\}.
\]

The right side specifies what the answer means before serialization. The
left side describes the proposed software path. This equality is an
obligation on the admitted domain, not a theorem established by two
programs returning parseable JSON. It includes a negative scientific
answer or a set with several components whenever that is the defined
output of \(F_Q\). A method that changes the target while improving
numerical stability has changed \(Q\); it has not merely implemented
\(E\).

There are three useful levels of equivalence. Exact structural
equivalence preserves a graph, variable map or expression tree up to
declared renaming. Numerical equivalence permits a quantified error in
the requested coordinates. Scientific equivalence additionally requires
the same population, intervention, sampling law and uncertainty
interpretation. A fixture showing agreement of two printed means
addresses only a small part of the second level. The design should state
which level it establishes and where further evidence is required.

An adapter also needs an explicit domain. Write \(\mathcal D_M\) for the
queries and inputs admitted by engine \(M\). If
\((D,Q)\notin\mathcal D_M\), the proper response is a typed unsupported
result that preserves the original request. Automatically dropping
weights, selecting one graph, clipping an interval or replacing a
temporal sample with independent rows does not extend the domain
honestly. A proposed approximation can still be useful, but its
transformed question and additional assumptions must be disclosed as a
different operation.

This view explains why a generic language bridge is insufficient. A
Python-to-R array conversion can preserve all floating-point values
while losing the fact that columns represent treatment arms in a
specific order. A Java graph object can preserve adjacency while
reversing an endpoint mark. A C++ routine can accurately optimize an
objective for a different weighting measure. The mathematical interface
must bind these meanings alongside the numeric payload, and its review
must include examples for which apparently harmless changes lead to
materially different bank decisions.

\hypertarget{population-measures-and-rate-coordinates}{%
\section{Population measures and rate
coordinates}\label{population-measures-and-rate-coordinates}}

Consider a bank's effect target

\[
 \tau_\mu=\int\{m_1(x)-m_0(x)\}\,d\mu(x),
\]

where \(m_a(x)\) is a conditional mean potential outcome under action
\(a\), and \(\mu\) is the declared distribution of customers or
exposures. The same response functions can imply different targets under
different measures. In an assumed two-segment example, the customer
shares are \(3/4\) and \(1/4\), the segment effects are USD 20 and USD
-40 per account, and relative balances are one and nine. The
customer-average effect is

\[
 \tfrac34(20)+\tfrac14(-40)=5.
\]

Balance weighting gives masses \(1/4\) and \(3/4\), and hence an effect
of USD -25. The sign reverses even though no within-segment estimate
changed. The example is stipulated arithmetic, not evidence about the
distribution of any bank's customers. Its purpose is to make the
weighting measure an indispensable part of the requested estimand.
\citep{exec13_interchange_examples}

A caller may intentionally prefer the balance-weighted target for a
treasury decision. That choice should create a new target identity and
explain which decision it serves. It cannot be inferred from a column
named \texttt{weight}, because a supplied weight might instead describe
unequal sampling probabilities, treatment assignment, transport to
another population or a numerical objective. These roles have different
identification and variance implications. Multiplying several weights
without preserving their definitions can conceal that distinction.

Treatment coordinates require similar care. Let \(r\) be a decimal
interest rate, so a basis-point change is \(10^{-4}\). If a linear
response derivative is \(\beta=\partial E(Y)/\partial r\), then the
derivative per basis point is \(10^{-4}\beta\). A hypothetical estimate
of USD 2 million per unit decimal rate becomes USD 200 per basis point;
an interval of USD 1--3 million becomes USD 100--300. Reporting the
original number under the new label would create a factor-of-10,000
error. The exact calculation is reproducible. It does not assert that a
constant derivative is an appropriate response model over an arbitrary
offer range. \citep{exec13_interchange_examples}

For a finite nonlinear rate change, the target is instead
\(E\{Y(r_1)-Y(r_0)\}\); rescaling a local derivative does not identify
that contrast. Nor does a per-month response become an annual causal
response by changing its unit label. Timing, compounding, feedback and
the intervention path may matter. Units should therefore describe the
mathematical coordinate actually estimated, while the horizon and
intervention path remain separate fields. A conversion that requires new
structural assumptions belongs in a scientific transformation, not a
formatting function.

\hypertarget{joint-uncertainty-across-coordinates}{%
\section{Joint uncertainty across
coordinates}\label{joint-uncertainty-across-coordinates}}

Suppose an admitted calculation returns \(\widehat\theta\) and
covariance \(\widehat\Sigma\) in the declared coordinates. Under a fixed
linear map \(\phi=L\theta+c\), the corresponding estimate and covariance
are

\[
 \widehat\phi=L\widehat\theta+c,
 \qquad \widehat\Sigma_\phi=L\widehat\Sigma L^\top.
\]

This covariance identity follows by expanding the centered vector
\(L(\widehat\theta-E\widehat\theta)\). It preserves whatever
justification the original covariance actually has; it supplies no new
consistency, normal approximation or finite-sample coverage argument. If
\(L\) is estimated, choosing to hold it fixed is a conditional analysis
and must be labeled. Joint estimation error cannot be eliminated by
storing the fitted \(L\) as a configuration value.

For an exact arithmetic illustration, let the first coordinate be a
gross revenue effect in USD and the second an operating-cost effect in
thousands of USD. Suppose the supplied covariance is

\[
 \Sigma=\begin{pmatrix}40000&20\\20&1/4\end{pmatrix},
 \qquad L=\begin{pmatrix}1/1000&0\\0&-1\end{pmatrix}.
\]

The new coordinates are gross benefit and negative operating cost, both
in thousands of USD. Their covariance is

\[
 L\Sigma L^\top=
 \begin{pmatrix}1/25&-1/50\\-1/50&1/4\end{pmatrix}.
\]

The negative cross-covariance matters when adding the two coordinates to
obtain net value: its variance is \(1/4\), whereas dropping that
covariance gives \(29/100\). A column-order error or a failure to
reverse the covariance sign changes the uncertainty even if both
marginal standard errors look reasonable. The matrices here are assumed,
positive definite reference objects whose conversion is checked by exact
arithmetic. They are not fitted bank uncertainty estimates.
\citep{exec13_interchange_examples}

Marginal intervals cannot substitute for a joint region. If \(C_1,C_2\)
each cover their own fixed coordinate with probability at least \(.95\),
the union bound gives only

\[
 P\{(\theta_1,\theta_2)\in C_1\times C_2\}\geq .90.
\]

Even with independent coverage events and exact marginal coverage, the
joint probability is \(.95^2=.9025\). A simultaneous 95\% statement
needs a construction that supports it; a transport cannot supply that
construction by changing the uncertainty label. The same issue arises if
an agent inspects several analyses and reports only the most favorable
one. Coverage for each fixed query is not automatically coverage for the
selected query. Preserving the full query-selection rule and any
simultaneous or selective argument is part of the workflow's scientific
identity.

These distinctions already appear in the local semantic validator:
confidence, posterior credibility and sensitivity have separate
meanings, a joint-confidence declaration requires construction evidence,
and the declared sampling law must match the context. Those checks
constrain what can be serialized. They do not independently prove the
submitted construction or establish Bayesian calibration.
\citep{exec13_interchange_runtime}

\hypertarget{a-confidence-set-is-not-a-pair-of-convenient-endpoints}{%
\section{A confidence set is not a pair of convenient
endpoints}\label{a-confidence-set-is-not-a-pair-of-convenient-endpoints}}

Weak identification supplies a practical reason to retain topology. For
a ratio target \(\theta=\mu_Y/\mu_X\), one can invert a test of
\(\mu_Y-\theta\mu_X=0\). With supplied approximately Gaussian mean
errors, write their variances as \(s_{YY},s_{XX}\) and covariance as
\(s_{XY}\). A residual-test set takes the form

\[
 (\bar Y-\theta\bar X)^2
 \leq c^2(s_{YY}-2\theta s_{XY}+\theta^2s_{XX}).
\]

Its interpretation depends on the stated Gaussian approximation,
covariance and critical value; algebraic inversion alone is not a
finite-sample theorem for an arbitrary dependent sample. Collecting
terms gives

\[
 a\theta^2+b\theta+d\leq0,
\]

where

\[
 a=\bar X^2-c^2s_{XX},\quad
 b=-2\bar X\bar Y+2c^2s_{XY},\quad
 d=\bar Y^2-c^2s_{YY}.
\]

When the denominator is weak, \(a\) may be negative. With two real roots
the solution then consists of exterior rays; with other coefficients it
may be the whole line or another topology. The point ratio can fail to
exist while the test-inversion set still has a precise meaning. For the
stipulated values \(\bar Y=2,\bar X=0,c=2,s_{YY}=s_{XY}=0,s_{XX}=1/4\),
the inequality is \(4\leq\theta^2\), giving

\[
 (-\infty,-2]\cup[2,\infty).
\]

This degenerate but positive-semidefinite reference covariance makes the
algebra transparent. It does not validate that covariance for a
data-generating process. It shows why a response object must allow a
valid set without a finite point ratio. The existing local
supplied-moment operation retains these distinctions under its declared
approximation.
\citep{exec13_interchange_examples,exec13_interchange_runtime}

Unknown, empty, whole-line and disconnected results answer different
questions. Unknown says that no set has been computed. Empty may signal
incompatible restrictions in a particular construction, and does not
mean an effect of zero. Whole-line uncertainty means every real value
remains allowed by that construction. An exterior-ray set excludes a
middle region despite being unbounded. A graph-class result may instead
be a set of distinct functionals whose numerical evaluation is not yet
available. None can be reduced to a nullable scalar without losing
information that could change a decision.

Nonlinear monetary transformations make the loss visible. Take the
assumed effect set \([-3,-2]\cup[2,3]\) and value map
\(V(\theta)=100(\theta^2-1)\). Its image is \([300,800]\). Applying the
same map to the convex hull \([-3,3]\) gives \([-100,800]\). Both are
outer possibilities for different retained information, but the latter
invents a negative-value possibility that the original set excludes. A
hull may be a disclosed conservative approximation; it must not be
described as exact interchange. The reverse mistake, returning only
transformed outer endpoints without considering interior extrema, can
omit valid values and lose coverage. \citep{exec13_interchange_examples}

\hypertarget{combining-statistical-measurement-and-numerical-uncertainty}{%
\section{Combining statistical, measurement and numerical
uncertainty}\label{combining-statistical-measurement-and-numerical-uncertainty}}

Suppose an effect set \(C\) contains the fixed target \(\theta\) with
probability at least \(1-\alpha_C\), and a set \(H\) contains a needed
valuation parameter \(\eta\) with probability at least \(1-\alpha_H\).
For a declared value function \(v\), the projection

\[
 W=\{v(t,h):t\in C,h\in H\}
\]

contains \(v(\theta,\eta)\) whenever both input events hold. Thus the
union bound yields coverage at least \(1-\alpha_C-\alpha_H\), without
independence. A valid joint region can be projected instead and may be
narrower because it retains dependence between \(\theta\) and \(\eta\).
Independently varying marginal ranges is generally an outer
approximation to that joint region. Fixing \(\eta\) at an estimate has
no such unconditional guarantee without an additional argument.

Numerical containment is a third obligation. Let \(\widetilde W\) be an
approximate computed set. If a certificate establishes

\[
 W\subseteq\widetilde W\oplus[-\varepsilon,\varepsilon],
\]

where \(\oplus\) denotes addition of sets, then the enlarged computed
set preserves the preceding coverage event. If the numerical certificate
itself can fail with probability \(\delta\), the resulting guarantee is
at least \(1-\alpha_C-\alpha_H-\delta\). Budgets of \(.02,.02,.01\) give
a 95\% lower guarantee. This arithmetic does not prove that a particular
optimizer supplies the assumed certificate; its error bound must come
from the numerical method. \citep{exec13_interchange_examples}

A grid of feasible points usually supplies an inner sample, not the
required outer enclosure. Missing a disconnected component or an
unbounded tail is therefore a scientific failure, even if all sampled
points satisfy the constraints. An optimizer's small objective change
also need not certify the set boundary. For a nonlinear map, a Lipschitz
constant can propagate an input approximation error only over a region
where the constant is valid; a bound obtained on a clipped plotting
window does not control unbounded tails outside it.

The currency of the numerical error should be the currency of the
decision. If a statistical lower value is USD 30 but an admitted
numerical approximation has error up to USD 40, the certified lower
value is USD -10. Reporting only the statistical endpoint can reverse a
claimed robust recommendation. Conversely, a deliberately coarse
computation may be adequate when its certified error cannot change the
decision. This makes numerical precision an explicit investment tradeoff
rather than an unexamined default tolerance.

\hypertarget{sampling-declarations-and-partial-scientific-answers}{%
\section{Sampling declarations and partial scientific
answers}\label{sampling-declarations-and-partial-scientific-answers}}

A shared data format does not create shared sampling assumptions.
Consider the generated one-date account model \(Y_i=\theta+C+E_i\),
where independent centered Gaussian components have
\(\operatorname{Var}(C)=.09\), \(\operatorname{Var}(E_i)=1\), and \(C\)
is shared by every account. Then

\[
 \operatorname{Var}(\bar Y)=.09+1/N.
\]

At \(N=10{,}000\), the correct variance is \(.0901\). Treating the rows
as independent with their individual variance \(1.09\) gives
\(.000109\), understating variance by \(90100/109\), or about 827 times.
The number of rows was preserved exactly; the scientific information was
not. This exact generated calculation illustrates why a panel or
calendar declaration must survive conversion and why an iid-only method
must reject it. It makes no claim about an observed bank covariance.
\citep{exec13_interchange_examples}

The same problem occurs with persistent linkage errors, cluster
treatment assignment and policies that condition on shared history. A
consumer that sees only row count and a floating-point estimate cannot
reconstruct the lost sampling unit. The local bounded back-door
reference explicitly restricts its sampling regime, while the proposed
symbolic identification adapter preserves the declaration without
supplying estimation uncertainty. Those are different services with
different obligations.
\citep{exec13_interchange_runtime,exec13_y0_contract}

Finally, the response should identify what has actually been
established. A source-inspected method may support a proposal decision.
An executed adapter may establish that particular conversions work. A
statistical region may be valid under declared scientific assumptions. A
banking action may still be undecided because its loss function or
information is insufficient. Preserving a partial answer makes the next
step precise: obtain a needed assumption, improve information, choose
another method or refine numerical containment. Automatically converting
every partial answer into a confident recommendation would obscure the
very gap the product is meant to explain.

\hypertarget{preserving-causal-meaning-in-a-python-identification-adapter}{%
\section{Preserving causal meaning in a Python identification
adapter}\label{preserving-causal-meaning-in-a-python-identification-adapter}}

\label{exec13:python-identification}

A common causal interface is useful only if a question keeps its meaning
when it passes into another library. An outreach decision can become a
different question if a transport drops a condition, converts an
uncertain graph into one directed graph, or treats a symbolic
probability formula as an estimated monetary return. The proposed Python
path therefore begins with a narrow service: identify an intervention
distribution in one supplied causal graph, preserve the assumptions that
make the formula relevant, and report an explicit limit when that
service cannot answer the question.

y0 is a plausible component for this service because its inspected
Python API separates a query, a mixed graph and a symbolic
identification expression. The source considered here is commit
\texttt{cc6644d3b470a1ab2a1eb9da4f0fdb0c843f5d3a}, dated 13 April 2026;
its project metadata declares development version \texttt{0.2.12-dev}.
The acquisition is later than the literature cutoff, but the reviewed
commit predates 7 October 2026. These identifiers describe inspected
source, not a tested deployment. The proposed decision is to integrate a
restricted identification path and verify its conversions before
admitting it to a prototype. Porting an entire identification engine
would create a second implementation to maintain before the first
interface has established its scientific meaning.
\citep{exec13_y0_api,exec13_y0_metadata}

\hypertarget{a-distributional-question-before-an-estimate}{%
\subsection{A distributional question before an
estimate}\label{a-distributional-question-before-an-estimate}}

Let \(A\) denote a bank action, \(Y\) an outcome and \(Z\) a collection
of explicitly named conditions. The requested object is

\[
 p(y\mid\operatorname{do}(a),z).
\]

It is distinct from an observed conditional probability, an average
difference between two actions, and a confidence interval for an
estimated difference. For a finite outcome support, a later evaluator
might form

\[
 \mu(a,z)=\sum_y y\,p(y\mid\operatorname{do}(a),z),
 \qquad \tau(z)=\mu(a_1,z)-\mu(a_0,z).
\]

That calculation requires the original outcome coordinates,
treatment/control order and units. If \(Y\) is a repayment indicator,
\(\tau\) is a probability difference. If \(Y\) is net cash received, it
has the declared currency and horizon. Multiplying the first by a loss
avoided per repayment is an additional valuation model; it is not an
operation that the identification library can infer from a variable
label.

The first adapter is deliberately limited to one scalar outcome and
treatment, optional conditions, one causal DAG or acyclic directed mixed
graph (ADMG), and declared finite discrete domains. This makes its
numerical references and interpretation concrete. It does not establish
an absence of wider functions elsewhere in y0. A graph representing
selection, transport between populations, cross-world counterfactuals,
an adaptive longitudinal regime or unresolved graph alternatives needs a
separately reviewed path. An account panel can still describe its
sampling law in the request, but a symbolic formula supplies no
independence assumption, standard error or calibration for that panel.

The scientific context should travel with the formula: target
population, action values, outcome horizon, observation and availability
clocks, variable domains, units, sign convention, sampling unit and
assumptions. A hash of that complete context identifies the request. A
changed population, a reversed action contrast or a different
availability clock produces a different identity even when the graph
picture is unchanged. The shared library, CLI and MCP operation
contracts already distinguish transport completion from scientific
disposition; the proposed adapter specializes that distinction for
identification. \citep{exec13_y0_contract}

A condition deserves particular care. Conditioning on a pretreatment
risk category and conditioning on a response to the action are different
targets. The interface should require the caller to state which is
intended, preserve the chronology, and avoid silently relabeling a
post-action variable as an adjustment covariate. Syntax can reject a
missing clock declaration. It cannot determine whether the bank's
claimed chronology or causal model is true.

\hypertarget{graph-and-coordinate-conversion}{%
\subsection{Graph and coordinate
conversion}\label{graph-and-coordinate-conversion}}

The graph object in the inspected source contains a directed component
and a component named \texttt{undirected}. In the causal mixed-graph
interpretation, the latter carries bidirected latent links: the source's
causal-graph importer maps the input type \texttt{bidirected} into that
field, and its latent expansion creates a shared hidden parent for each
such edge. It is therefore incorrect to put an unresolved CPDAG edge
into this component merely because both are stored by an undirected
graph data structure. \citep{exec13_y0_graph}

For the admitted path, a tail--arrow pair becomes a directed edge, and
an arrow--arrow pair becomes an ADMG bidirected edge. Circles,
tail--tail edges and graph collections receive an unsupported
disposition that retains the original representation. A directed edge
and a bidirected edge may coexist on the same pair of variables; erasing
one would change the model. The directed component must be acyclic, with
all endpoints in the declared node set. These checks belong in the
adapter: the inspected graph initializer sets metadata flags and does
not itself establish acyclicity.

Every variable, including an isolated one, must be passed explicitly.
Constructing a graph solely from edges can lose an outcome or an
independent condition that has no incident edge. The adapter should
compare the node and typed-edge sets before and after conversion and use
a deterministic valid topological ordering. The purpose is
reproducibility and semantic preservation, not the selection of a
statistically preferred causal ordering.

Names are another coordinate system. External identifiers may contain
spaces, non-ASCII text or labels that collide after informal
sanitization. The proposed conversion assigns opaque internal symbols
through a bijection and keeps a sidecar containing the external
identifier, displayed label, domain order, time coordinate and units. It
never guesses temporal order from a name such as \texttt{lag1}, and it
never converts category labels into numeric doses without a declared
mapping. A graph-set request remains a graph-set request even when
unsupported. An uncomputed effect set remains unknown; it is not
represented by a convenient point estimate or an empty region.

\hypertarget{why-conditional-identification-needs-a-restricted-route}{%
\subsection{Why conditional identification needs a restricted
route}\label{why-conditional-identification-needs-a-restricted-route}}

The high-level API calls ID when its conditions field is empty and IDC
otherwise. By default it converts the specific \texttt{Unidentifiable}
exception to \texttt{None}; with \texttt{strict=True} it preserves the
exception. Other exception types are not caught by that branch. The
adapter should preserve this distinction, because a malformed graph, a
runtime failure and nonidentification are different outcomes. The
inspected ID recursion supplies node/district arguments when raising
\texttt{Unidentifiable}, but those arguments should not be presented as
an independently verified graphical certificate.
\citep{exec13_y0_api,exec13_y0_id}

Two additional source findings restrict the proposed conditional route.
First, \texttt{Query.expression} builds a conditioned distribution when
conditions exist, using an \texttt{elif} branch for interventions. When
both are present, that convenience property does not add the
intervention. This is a serialization hazard, not evidence that the
separate IDC algorithm discards its treatment field:
\texttt{Identification} retains outcomes, treatments and conditions
separately. The adapter should serialize its original query triple
independently of that convenience property. \citep{exec13_y0_utils}

Second, IDC tests whether an observation may be exchanged for an action
using a separation routine. The inspected routine takes a
directed-ancestor subgraph, connects directed co-parents, flattens the
remaining edges and removes conditioning vertices. For a DAG this is a
familiar construction. The same steps do not generally implement the
collider behavior of bidirected paths. The source-level discrepancy can
be demonstrated without relying on a package's own numerical output.
\citep{exec13_y0_idc,exec13_y0_rule2,exec13_y0_separation,exec13_y0_graph}

Consider the ADMG

\[
 Y\leftrightarrow B\leftrightarrow Z,
\]

with an unrelated action \(X\). The middle vertex \(B\) is a collider.
Conditioning on it can connect \(Y\) and \(Z\); treating it as an
ordinary undirected intermediate vertex and deleting it instead
disconnects them. In the rule-2 check for exchanging \(Z\), removal of
incoming edges to isolated \(X\) and outgoing edges from \(Z\) does not
change this path. The inspected directed co-parent rule adds no
\(Y\)--\(Z\) link because there are no directed edges. Its subsequent
deletion of \(B\) therefore predicts separation where the causal mixed
graph does not.

An exact generated law makes the consequence explicit. Let \(U,V\) be
independent fair bits, let \(Y=U\), \(Z=V\), and let \(B\) report
whether \(U=V\), correctly with probability \(.9\). The observed law of
\((Y,B,Z)\) has positive probability in all eight cells. Direct
enumeration gives

\[
 P(Y=1\mid B=1,Z=1)=\frac9{10},
 \qquad
 P(Y=1\mid B=1,\operatorname{do}(Z=1))=\frac12.
\]

An intervention on \(Z\) overwrites that measurement but does not change
\(U,V\) or the mechanism generating \(B\). Exchanging the observation
for the action would therefore change the requested probability. This is
an exact reference for the scientific meaning of the graph and a trace
of the inspected source path; it is not a report of an executed y0
failure. \citep{exec13_y0_oracles}

The proposed adapter can avoid that native conditional path, subject to the internal node-preservation qualification in Section~\ref{exec13:isolated-condition}. Its mathematical route first applies ID to the joint query \(p(y,z\mid\operatorname{do}(a))\). If it
obtains an observational functional \(N(y,z,a)\), it forms

\[
 \frac{N(y,z,a)}{\sum_{y'}N(y',z,a)},
\]

on coordinates for which the denominator is positive. This is a direct
probability identity. It neither uses the questionable separation step
nor relies on the query-expression convenience property. The route is
conservative: failure to identify the joint distribution does not prove
failure to identify the conditional target. Such a failure must return
``unsupported by the admitted conditional route,'' leaving the original
question intact. Re-enabling native IDC requires a corrected or
otherwise justified path and an actual engine comparison against
independent collider references.

\hypertarget{a-front-door-reference-with-a-banking-interpretation}{%
\subsection{A front-door reference with a banking
interpretation}\label{a-front-door-reference-with-a-banking-interpretation}}

A second exact example supplies a positive identification reference. Let
\(A\) be assistance, \(M\) completion of the assisted process and \(Y\)
successful repayment. In the generated example an unobserved \(U\)
affects both assistance and repayment. The assumed graph has
\(A\to M\to Y\) and \(A\leftrightarrow Y\), with no direct effect of
assistance on repayment outside \(M\). This is a strong exclusion
assumption. In a real service workflow an agent's advice, fees or other
treatment channels could violate it, and observed associations alone do
not establish it.

Specify the finite law by

\[
 \begin{aligned}
 P(U=1)&=\tfrac12,\\
 P(A=1\mid U=u)&=\tfrac15+\tfrac35u,\\
 P(M=1\mid A=a)&=\tfrac14+\tfrac12a,\\
 P(Y=1\mid M=m,U=u)&=\tfrac1{10}+\tfrac12m+\tfrac15u.
 \end{aligned}
\]

Independent rational enumeration first constructs the observed joint
table, hiding \(U\), and evaluates

\[
 \sum_m p(m\mid a)\sum_{a'}E(Y\mid m,a')p(a').
\]

It separately evaluates the intervention law from the specified
mechanisms. Both routes produce \(13/40=.325\) under \(a=0\) and
\(23/40=.575\) under \(a=1\), for a difference of \(1/4\). The naive
observed difference is \(37/100\). The exact calculation therefore
distinguishes a successful symbolic formula from a shortcut that merely
compares observed groups. No fitted sample, randomized qualification
study or engine runtime is involved. \citep{exec13_y0_oracles}

If a proposal scenario assigns an assumed USD 1,000 benefit to each
additional successful repayment, the gross expected difference is USD
250 per account, before assistance costs and any other consequences.
That number is planning arithmetic under the declared law and value
assumption. It supplies neither observed bank ROI nor evidence that the
front-door graph describes a bank.

The inner treatment variable \(a'\) reveals an important serialization
requirement. It is a dummy summation coordinate, while the outer \(a\)
is the requested intervention value. A global string replacement could
capture the dummy variable and change the formula. The proposed result
format uses a typed expression tree with lexical bindings for
marginalized variables. It preserves probability factors, products, sums
and ratios, alpha-renames bound symbols when needed, and maps free
coordinates back to the original question. The inspected symbolic
classes expose those structures; rendered LaTeX or the library's
Python-like display string is suitable for presentation, not as an
executable interchange format. \citep{exec13_y0_dsl}

An independent negative reference is equally necessary. In a graph
allowing \(A\to Y\) and \(A\leftrightarrow Y\), take a fair hidden
\(U\), let \(A=U\), and let an independent bit \(E\) equal one with
probability \(.2\). The mechanisms \(Y=A\mathbin{\mathrm{xor}}E\) and
\(Y=U\mathbin{\mathrm{xor}}E\) give the same strictly positive observed
law of \((A,Y)\). Under intervention their effects are respectively
\(.6\) and zero. Both belong to the model allowing those edges; no
faithfulness restriction is imposed. A library cannot choose between the
effects from their identical observed distribution. A successful adapter
retains this nonidentification rather than translating every completed
request into a finite effect. \citep{exec13_y0_oracles}

\hypertarget{responses-deployment-and-the-next-prototype-decision}{%
\subsection{Responses, deployment and the next prototype
decision}\label{responses-deployment-and-the-next-prototype-decision}}

A symbolic identification result should contain the original query,
graph and context identities, engine revision, typed functional, support
conditions and source provenance. Its numeric estimate, standard error
and confidence region remain absent. When the unconditional ID path
reports nonidentification in an admitted model, the response should
identify that origin and retain its diagnostic without exaggerating its
verification. When the conditional joint route fails, the response is
unsupported. An invalid graph is rejected; a missing engine, internal
failure, timeout or cancellation is unavailable. These distinctions
determine whether the next action is to repair an input, add scientific
information, select a different method or provision a runtime.

Library, CLI and MCP surfaces should call the same core implementation
and return the same context and result identities. The symbolic service
requires graph and question artifacts, not bank account rows. A worker
can enforce wall-time, memory and output budgets while keeping engine
diagnostics separate from scientific conclusions. Cancellation must
preserve a completed result if publication has already occurred; an
interrupted partial expression cannot become an identified functional.
These are proposed implementation requirements, not performance
guarantees supplied by the inspected y0 API.

The pinned project metadata requires Python 3.11 or later and lists
twelve direct runtime dependencies, including NetworkX, SciPy, pandas,
scikit-learn and statsmodels. The build backend is \texttt{uv\_build} in
its declared version range. Its optional R extra uses \texttt{rpy2};
upstream test groups include additional packages that are not all
runtime requirements of the proposed pure-Python path. The inspected
tree has no root lockfile establishing resolved transitive versions. A
prototype admission therefore needs a frozen interpreter/platform and
dependency manifest, saved distribution hashes, actual build/import logs
and license notices. Source inspection alone cannot promise
compatibility with whatever environment happens to be installed.
\citep{exec13_y0_metadata,exec13_y0_tree}

The saved BSD three-clause license permits redistribution under its
notice, disclaimer and nonendorsement conditions. That supports a
possible integration route, while leaving dependency rights and actual
deployment behavior to be checked. It supplies no scientific warranty. A
port may become justified if a later measurement shows that the Python
environment or latency is a decisive constraint; there is no such
measured advantage for porting in the present evidence.
\citep{exec13_y0_license}

The prototype's first engine trials should compare the actual pinned
runtime with the independent front-door, bow and collider references,
including retained unsupported requests. Conversion checks should
exercise isolated variables, edge-type distinctions, coordinate
permutations, dummy-variable binding, zero conditional support and
distinct failures. A successful trial would support the named adapter in
its declared regime. It would not establish discovery accuracy, graph
truth, estimation calibration or institutional acceptance. The value of
this restricted design is that each later claim has a specific test and
a specific scientific object that must survive the interface. The next study tests that requirement inside graph modification, where a correctly converted isolated condition can still be lost; the later comparisons extend the same preservation question to custody, recovery and other engines.
\label{exec13:python-identification-end}

\label{exec13:interchange-end}
